%%%%%
%
\section{Humidit\'e absolue issue d'un couplage entre le Lyman-$\alpha$ 
et le TD}

\itemstyle{Op\'erateur~:} {\sf ha\_lyman\_td}
\index{ha\_lyman\_td}

%%%%%OBJET%%%%
\itemstyle{Objet~:}
Calcul de l'humidit\'e absolue rapide \`a partir de la 
mesure du  Lyman-$\alpha$ coupl\'ee avec la mesure du capteur
\`a point de ros\'ee.

%%%%%ARGUMENTS%%%%
\bigskip
%%%%%ENTREES%%%%
\begin{tabarg}
\multicolumn{3}{l}{\titretab {Entr\'ees}}\\
\hline
$T_d$ & Vecteur, {\tt float} &  Temp\'erature du point de ros\'ee
($K$ ou $^\circ C$) \\
$V_{pel}$ & Vecteur, {\tt float} & Tension Peltier du capteur
\`a point de ros\'ee  ($V$)\\
$V_{lym}$ & Vecteur, {\tt float} & Tension de sortie du
capteur Lyman-$\alpha$ ($V$)\\
$T_s$ & Vecteur, {\tt float} & Temp\'erature de l'air 
($K$ ou $^\circ C$)\\
$P_s$ & Vecteur, {\tt float} & Pression statique ($hPa$) \\
$V_p$ & Vecteur, {\tt float} & Vitesse propre ($m/s$)\\
$f_e$ & Coeff., {\tt float} & Fr\'equence d'\'echantillonnage ($Hz$)\\
$V_{pelI}$ & Coeff., {\tt float} & Borne inf\'erieure de validit\'e
de la tension Peltier ($V$)\\
$V_{pelS}$ & Coeff., {\tt float} & Borne sup\'erieure de validit\'e
de la tension Peltier ($V$)\\
\end{tabarg}

\medskip
%%%%%SORTIES%%%%
\begin{tabarg}
\multicolumn{3}{l}{\titretab {Sorties}}\\
\hline
$H_a$ & Vecteur, {\tt float} & Humidit\'e absolue ($g/m^3$)\\
\end{tabarg}


%%%%%FORMULE%%%%
\itemstyle{Formule}

%%ha_lyman_td
Dans un premier temps, l'algorithme trie les valeurs issues du
capteur \`a point de ros\'ee (TD) pour garder
uniquement les points o\`u le capteur n'est pas en phase de recherche.
La s\'election se fait suivant
deux crit\`eres~:\\
\indent - valeur de la tension de l'\'el\'ement Peltier proche de
la valeur \`a l'\'equilibre (compris entre $V_{pelI}$ et $V_{pelS}$)\\
\indent - $Td \leq Ts$ (suppression des d\'epassements)\\
\\
La serie d'humidit\'e absolue de r\'ef\'erence est calcul\'ee \`a partir
de la s\'erie des valeurs correctes du TD.\\
A chaque pas de temps,
une r\'egression lin\'eaire est r\'ealis\'ee entre, les
N points d'humidit\'e absolue de r\'ef\'erence et, les N points
correspondants
du Lyman-$\alpha$, encadrant la valeur \`a calibrer.
La validation du r\'esultat de cette r\'egression est soumise \`a
deux crit\`eres~:\\
\indent - \'ecart-type de l'humidit\'e absolue sup\'erieur \`a 0.5 $g/m^3$
(dynamique du signal suffisante)\\
\indent - coefficient de corr\'elation sup\'erieur \`a 0.8 (qualit\'e de
la r\'egression)\\



%%%%COMMENTAIRE%%%
\itemstyle{Commentaire~:}

Cet op\'erateur est un essai de couplage entre 
le Lyman-$\alpha$ et le capteur \`a point de
ros\'ee. Il est tr\`es 
difficile d'ajuster les crit\`eres de 
validit\'e de la r\'egression pour le rendre
robuste.
Donc, cet op\'erateur doit \^etre utilis\'e
avec prudence.

%%%%%AUTEUR%%%%
\itemstyle{Auteur~:}
CNRM/GMEI/TRAMM

%%%%%REFERENCE%%%%
\itemstyle{R\'ef\'erences~:}
